\documentclass[11pt,a4paper]{article}

% --- Mathematics (must precede unicode-math) ---
\usepackage{amsmath}
\usepackage{mathtools}

% --- Fonts (lualatex) ---
\usepackage{fontspec}
\usepackage{unicode-math}
\setmathfont{KpMath-Regular.otf}
\setmathfont{KpMath-Bold.otf}[version=bold]

% --- Microtypography (after all font setup) ---
\usepackage{microtype}

% --- Colors & page layout ---
\usepackage[dvipsnames]{xcolor}
\usepackage[margin=25mm, top=30mm, bottom=30mm]{geometry}

% --- Headers & footers ---
\usepackage{fancyhdr}
\pagestyle{fancy}
\fancyhf{}
\fancyhead[L]{\small\itshape P--SV--SH Optical Theorem}
\fancyhead[R]{\small\itshape Nestor (2026)}
\fancyfoot[C]{\thepage}
\renewcommand{\headrulewidth}{0.4pt}
\renewcommand{\footrulewidth}{0pt}

% --- Section numbering ---
\setcounter{secnumdepth}{3}

% --- Tables ---
\usepackage{booktabs}
\usepackage{array}

% --- Captions ---
\usepackage[font=small, labelfont=bf, labelsep=period]{caption}

% --- Spacing ---
\usepackage{setspace}
\setstretch{1.15}
\setlength{\parskip}{0.4em}
\setlength{\emergencystretch}{3em}

% --- Bibliography ---
\usepackage[round,authoryear]{natbib}
\bibliographystyle{plainnat}

% --- Hyperlinks (last before macros) ---
\usepackage[colorlinks=true, linkcolor=blue!60!black,
            citecolor=blue!60!black, urlcolor=blue!60!black]{hyperref}

% --- Macros (house set) ---
\newcommand{\bI}{\mathbf{I}}
\newcommand{\bx}{\mathbf{x}}
\newcommand{\dd}{\mathrm{d}}
\newcommand{\bT}{\mathbf{T}}

% --- New macros for this document ---
\newcommand{\bS}{\mathbf{S}}
\newcommand{\bu}{\mathbf{u}}
\newcommand{\bA}{\mathbf{A}}
\newcommand{\bJ}{\mathbf{J}}
\newcommand{\bd}{\hat{\mathbf{d}}}
\newcommand{\bp}{\hat{\mathbf{p}}}
\newcommand{\xh}{\hat{\mathbf{x}}}
\newcommand{\thh}{\hat{\boldsymbol{\theta}}}
\newcommand{\phh}{\hat{\boldsymbol{\phi}}}
\newcommand{\bsig}{\boldsymbol{\sigma}}
\newcommand{\herm}{^{\dagger}}
\newcommand{\sext}{\sigma_{\mathrm{ext}}}
\newcommand{\ssc}{\sigma_{\mathrm{sc}}}
\newcommand{\sabs}{\sigma_{\mathrm{abs}}}
\newcommand{\kp}{k_{P}}
\newcommand{\ks}{k_{S}}
\DeclareMathOperator{\Imag}{Im}
\DeclareMathOperator{\Real}{Re}

\begin{document}

% --- Title block ---
\thispagestyle{plain}
\begin{center}
\vspace*{1em}
{\LARGE\bfseries P--SV--SH Optical Theorem:\\[0.3em]
Forward Amplitude, Flux Weights and Extinction}\\[1.8em]
{\large T.\,M.\ Nestor}\\[0.5em]
{\itshape Research School of Earth Sciences, Australian National University}\\[2em]
\end{center}

\noindent\textbf{Status.} Explanatory note. The results of
\S\S\ref{sec:setup}--\ref{sec:theorem} are textbook material
\citep{Varatharajulu1977,PaoMow1973,Newton1982}. In this project the
$S$-incidence form \eqref{eq:ot-s} has been checked numerically for the
lossless sphere to $10^{-8}$ (\texttt{FirstOrderContrastOperator.tex},
toroidal-term paragraph). The $P$-incidence form \eqref{eq:ot-p} is
checked for the exact Mie sphere and for the $T_{27}$ cube
(\S\ref{sec:sign}). The cube symmetry statements of
\S\ref{sec:forward-matrix} and the T-matrix form of \S\ref{sec:tmatrix}
have no independent check here and are \emph{unvalidated} in the sense of
the project's engineering standards.

% =============================================================================
\section{Setting and far field}
\label{sec:setup}
% =============================================================================

A bounded scatterer $V$ is embedded in a homogeneous isotropic elastic
background with density $\rho$ and wave speeds $\alpha$ (P) and $\beta$
(S). The time dependence is $e^{-i\omega t}$ throughout, and
\begin{equation}
  \kp = \omega/\alpha, \qquad \ks = \omega/\beta, \qquad \kp < \ks .
  \label{eq:k}
\end{equation}
Outside $V$ the scattered displacement satisfies the Kupradze radiation
condition \citep{PaoVaratharajulu1976}. Far from $V$ it separates into
an outgoing longitudinal wave and an outgoing transverse wave
\citep{PaoMow1973,AkiRichards2002}:
\begin{equation}
  \bu^{s}(\bx) \sim
  \bA_P(\xh)\,\frac{e^{i\kp r}}{r}
  + \bA_S(\xh)\,\frac{e^{i\ks r}}{r},
  \qquad r = |\bx| \to \infty,
  \label{eq:farfield}
\end{equation}
where $\bA_P \parallel \xh$ and $\bA_S \perp \xh$. The transverse
amplitude splits in the spherical triad into
\begin{equation}
  \bA_S = A_{SV}\,\thh + A_{SH}\,\phh ,
  \label{eq:svsh}
\end{equation}
so the far field carries three channels, P, SV and SH. The amplitudes
$\bA_P$ and $\bA_S$ are displacement amplitudes with dimension of
length, per unit incident amplitude.

Two incident waves are considered, each of unit displacement amplitude:
\begin{align}
  \text{P incidence:}\quad & \bu^{i} = \bd\, e^{i\kp \bd\cdot\bx},
  \label{eq:inc-p}\\
  \text{S incidence:}\quad & \bu^{i} = \bp\, e^{i\ks \bd\cdot\bx},
  \qquad \bp\cdot\bd = 0,\ |\bp| = 1 .
  \label{eq:inc-s}
\end{align}
For S incidence $\bp$ may be the SV or SH direction, or complex for
elliptical polarisation.

% =============================================================================
\section{Energy flux and channel weights}
\label{sec:flux}
% =============================================================================

For time-harmonic motion the time-averaged energy flux vector is
\begin{equation}
  \langle \bJ \rangle = \tfrac{\omega}{2}\,
    \Imag\bigl(\bsig\cdot\bar{\bu}\bigr),
  \label{eq:flux}
\end{equation}
with $\bsig$ the stress tensor \citep{Achenbach1973,AkiRichards2002}.
For a plane wave of amplitude $U$ travelling at speed $c$ this reduces
to $\tfrac12\rho\omega^2 c\,|U|^2$ along the propagation direction.
Applied to each term of \eqref{eq:farfield}, the scattered power per
unit solid angle is
\begin{equation}
  \frac{\dd P_{s}}{\dd\Omega}
  = \tfrac12\rho\omega^{2}
    \bigl(\alpha\,|\bA_P|^{2} + \beta\,|\bA_S|^{2}\bigr).
  \label{eq:dP}
\end{equation}
The P--S cross term drops out of \eqref{eq:dP}: its phase
$e^{i(\ks-\kp)r}$ oscillates in $r$, and its polarisations are
orthogonal.

Cross sections are powers divided by the incident flux, which is
$\tfrac12\rho\omega^2 \alpha$ for \eqref{eq:inc-p} and
$\tfrac12\rho\omega^2 \beta$ for \eqref{eq:inc-s}. Each scattered
channel is therefore weighted by the ratio of its wave speed to the
incident speed \citep{WuAki1985,SatoFehlerMaeda2012}:
\begin{align}
  \text{P incidence:}\quad
  \ssc &= \oint |\bA_P|^{2}\,\dd\Omega
        + \frac{\beta}{\alpha}\oint |\bA_S|^{2}\,\dd\Omega ,
  \label{eq:ssc-p}\\
  \text{S incidence:}\quad
  \ssc &= \frac{\alpha}{\beta}\oint |\bA_P|^{2}\,\dd\Omega
        + \oint |\bA_S|^{2}\,\dd\Omega .
  \label{eq:ssc-s}
\end{align}
Since $\beta/\alpha = \kp/\ks$, the weights are ratios of wavenumbers
too. The S term in \eqref{eq:ssc-p} and \eqref{eq:ssc-s} is the sum of
its SV and SH parts, $|\bA_S|^2 = |A_{SV}|^2 + |A_{SH}|^2$.

% =============================================================================
\section{The theorem}
\label{sec:theorem}
% =============================================================================

Write the total field as $\bu = \bu^i + \bu^s$ and integrate the flux
\eqref{eq:flux} over a large sphere $S_R$. The net outward flux is
$-P_{\mathrm{abs}}$, where $P_{\mathrm{abs}}$ is the power absorbed in $V$.
The incident term alone integrates to zero. The scattered term gives
$P_s$ through \eqref{eq:dP}. What remains is the interference term,
\begin{equation}
  P_{\mathrm{ext}}
  = -\tfrac{\omega}{2}\oint_{S_R}
     \Imag\bigl(\bsig^{i}\cdot\bar{\bu}^{s}
               + \bsig^{s}\cdot\bar{\bu}^{i}\bigr)\cdot\xh\;\dd S ,
  \label{eq:pext}
\end{equation}
with the balance
\begin{equation}
  P_{\mathrm{ext}} = P_{s} + P_{\mathrm{abs}}, \qquad
  \sext = \ssc + \sabs .
  \label{eq:balance}
\end{equation}
The integrand of \eqref{eq:pext} carries the phase
$e^{ik(r - \bd\cdot\bx)}$, which oscillates rapidly on $S_R$ as
$R\to\infty$. Stationary phase \citep{MorseFeshbach1953} confines
the integral to the forward point $\xh = \bd$, where the phase is
stationary and independent of $R$. The backward point contributes a term
proportional to $e^{2ikR}$. Equation \eqref{eq:balance} must hold for every
$R$, so that term vanishes. This is the same argument as in the scalar
and electromagnetic cases \citep{Newton1976,BohrenHuffman1983}.

The elastic case adds one step. The incident wave can interfere only
with the scattered wave of the \emph{same} type:
\begin{itemize}
\item \textbf{Different wavenumbers.} Incident P interferes with
  scattered S with the phase $e^{i(\ks - \kp)R}$ at the forward point,
  and the reverse pairing is the same. This phase is never stationary in
  $R$, so the term vanishes by the argument above.
\item \textbf{Orthogonal polarisations.} At $\xh = \bd$ the scattered S
  wave is transverse, $\bA_S(\bd)\perp\bd$, while an incident P wave is
  polarised along $\bd$. Conversely, a scattered P wave is polarised
  along $\bd$, while an incident S wave is polarised along $\bp\perp\bd$.
  The leading-order cross flux is proportional to the polarisation
  product and is zero.
\end{itemize}
Evaluating the same-mode forward term gives the elastic optical theorem,
or forward-amplitude theorem \citep{Varatharajulu1977}:
\begin{align}
  \text{P incidence:}\quad
  \sext &= \frac{4\pi}{\kp}\,
           \Imag\bigl[\bd\cdot\bA_P(\bd)\bigr],
  \label{eq:ot-p}\\
  \text{S incidence:}\quad
  \sext &= \frac{4\pi}{\ks}\,
           \Imag\bigl[\bar{\bp}\cdot\bA_S(\bd)\bigr].
  \label{eq:ot-s}
\end{align}
For a lossless scatterer $\sabs = 0$, so the imaginary part of one
forward amplitude sets the total power scattered into all three
channels P, SV and SH, with the weights of \eqref{eq:ssc-p} and
\eqref{eq:ssc-s}. Mode conversion is fully accounted for, but only on
the right of \eqref{eq:balance}. It never appears in $\sext$.

\begin{table}[h]
\centering
\caption{Contents of the two sides of \eqref{eq:balance} for each
incident type. ``Forward'' means $\xh = \bd$.}
\label{tab:summary}
\begin{tabular}{l >{\raggedright}p{0.33\textwidth} >{\raggedright\arraybackslash}p{0.38\textwidth}}
\toprule
Incident & $\sext$ uses & $\ssc$ uses \\
\midrule
P & forward P amplitude, component along $\bd$
  & $|\bA_P|^2$ with weight 1; $|A_{SV}|^2 + |A_{SH}|^2$ with weight $\beta/\alpha$ \\[0.3em]
SV & forward S amplitude, $\thh$ component
  & $|\bA_P|^2$ with weight $\alpha/\beta$; $|A_{SV}|^2 + |A_{SH}|^2$ with weight 1 \\[0.3em]
SH & forward S amplitude, $\phh$ component
  & as for SV \\
\bottomrule
\end{tabular}
\end{table}

% =============================================================================
\section{The forward S matrix and polarisation}
\label{sec:forward-matrix}
% =============================================================================

For S incidence the forward response is linear in $\bp$. In the
$(\thh,\phh)$ basis at $\xh = \bd$ it is a $2\times2$ matrix
$\mathsf{F}(\bd)$,
\begin{equation}
  \bA_S(\bd) = \mathsf{F}(\bd)\,\bp, \qquad
  \sext(\bp) = \frac{4\pi}{\ks}\,\Imag\bigl[\bp\herm\mathsf{F}(\bd)\,\bp\bigr].
  \label{eq:fmatrix}
\end{equation}
Off-diagonal elements of $\mathsf{F}$ (forward SV$\to$SH conversion)
enter $\sext$ only through the anti-Hermitian part of $\mathsf{F}$. For
linear SV or SH polarisation, the extinction takes the corresponding
diagonal element. Forward cross-polarised power still counts in
$\ssc$ through $|\bA_S|^2$.

The structure of $\mathsf{F}$ follows from the point symmetry of the
scatterer about the axis $\bd$. It must commute with every symmetry
operation that fixes $\bd$.
\begin{itemize}
\item \textbf{Sphere.} Full axial symmetry, so $\mathsf{F}$ is scalar,
  and SV and SH extinction are equal for any $\bd$.
\item \textbf{Cube, $\bd$ along $[100]$ or $[111]$.} The symmetry groups
  are $C_{4v}$ and $C_{3v}$. Each contains a rotation of order at least
  three and a mirror containing the axis, and this forces $\mathsf{F}$ to
  be scalar. Extinction is independent of polarisation.
\item \textbf{Cube, $\bd$ along $[110]$.} The group is $C_{2v}$, with
  mirrors $(001)$ and $(1\bar{1}0)$. In the basis aligned with them,
  $\mathsf{F}$ is diagonal with two different eigenvalues in general.
  The two linear polarisations have different extinction, so a cube is
  weakly dichroic, and weakly birefringent, along face diagonals.
\item \textbf{Cube, general $\bd$.} No symmetry. $\mathsf{F}$ is a
  general $2\times2$ matrix, and the eigenpolarisations need not be
  linear.
\end{itemize}
The same argument applied to P incidence gives nothing new, since
$\bd\cdot\bA_P(\bd)$ is a scalar for every scatterer.

% =============================================================================
\section{T-matrix form: flux-normalised unitarity}
\label{sec:tmatrix}
% =============================================================================

In a spherical-wave basis the regular and outgoing vector wavefunctions
split into a longitudinal family $\mathbf{L}$ (P), and two transverse
families, $\mathbf{M}$ (toroidal, SH-type) and $\mathbf{N}$ (spheroidal,
SV-type) \citep{PaoMow1973,Waterman1976}. The T-matrix maps regular
incident coefficients to outgoing scattered coefficients. The scattering
matrix
\begin{equation}
  \bS = \bI + 2\bT
  \label{eq:s-from-t}
\end{equation}
maps incoming to outgoing waves \citep{VaratharajuluPao1976,Newton1982}.

If each basis function carries unit energy flux, the energy balance
\eqref{eq:balance} becomes, for all incident vectors,
\begin{equation}
  \bI - \bS\herm\bS \succeq 0,
  \qquad \text{with equality when the scatterer is lossless.}
  \label{eq:unitary}
\end{equation}
Substituting \eqref{eq:s-from-t} gives, in the lossless case,
\begin{equation}
  \bT + \bT\herm = -2\,\bT\herm\bT,
  \qquad
  -\Real T_{\nu\nu} = \sum_{\mu} |T_{\mu\nu}|^{2},
  \label{eq:t-unitary}
\end{equation}
where $\mu$ runs over every outgoing channel in the $\mathbf{L}$,
$\mathbf{M}$ and $\mathbf{N}$ families. Equation \eqref{eq:t-unitary}
is the optical theorem written one partial wave at a time: the left
side is extinction, the right side is scattering. Summing over the
partial waves of a plane wave and evaluating the result in the forward
direction recovers \eqref{eq:ot-p} and \eqref{eq:ot-s}
\citep{Waterman1976,Varatharajulu1977}.

\paragraph{Normalisation is essential.} Standard vector wavefunctions
are not flux-normalised. The $\mathbf{L}$ family is built on $\kp$ and
the $\mathbf{M}$, $\mathbf{N}$ families on $\ks$, so a raw T-matrix
satisfies \eqref{eq:t-unitary} only with channel weights. These weights
are the same speed ratios as in \eqref{eq:ssc-p} and \eqref{eq:ssc-s}.
Without them, the P--S blocks appear to violate energy conservation by
factors of $\beta/\alpha$ or $\alpha/\beta$. The same rescaling makes
the reciprocity relation a plain transpose, $\bS = \bS^{\mathsf T}$
\citep{Varatharajulu1977,Pao1978}. The analogous statement for
plane-wave interface coefficients is in \texttt{psvsh\_unitarity.tex}.

% =============================================================================
\section{Coherent attenuation}
\label{sec:foldy}
% =============================================================================

The theorem is what connects a single-scatterer T-matrix to the
attenuation of the coherent wave in a random distribution. For number
density $n$ in the Foldy approximation \citep{Foldy1945,Lax1951,Martin2006},
the coherent P wave has effective wavenumber
\begin{equation}
  K_P^{2} \approx \kp^{2} + 4\pi n\,\bd\cdot\bA_P(\bd),
  \label{eq:foldy}
\end{equation}
so to first order in $n$,
\begin{equation}
  \Imag K_P \approx \frac{2\pi n}{\kp}\,\Imag\bigl[\bd\cdot\bA_P(\bd)\bigr]
  = \tfrac12\, n\,\sext .
  \label{eq:atten}
\end{equation}
The coherent amplitude therefore decays at half the extinction per unit
length, and the coherent intensity decays at the full rate
$n\sext$ \citep{SatoFehlerMaeda2012}. The S case is identical,
with $\mathsf{F}(\bd)$ taking the place of the scalar forward amplitude.
Along a direction where $\mathsf{F}$ is not scalar, such as $[110]$ for
aligned cubes, the coherent S wave splits into two eigenpolarisations
with different $K_S$.

% =============================================================================
\section{Uses as a numerical check}
\label{sec:check}
% =============================================================================

For a lossless scatterer, the ratio $\sext/\ssc$ is unity with no free
parameters. The two sides are computed from different data: $\sext$
from one forward amplitude, $\ssc$ from an angular quadrature over all
three channels. The ratio therefore checks any T-matrix independently
of how the T-matrix was obtained. Its main uses are these.
\begin{itemize}
\item \textbf{Truncation.} A deficit that shrinks with increasing
  $n_{\max}$ measures multipole truncation.
\item \textbf{Mode-conversion weights.} Dropping the flux weight in
  \eqref{eq:ssc-p} or \eqref{eq:ssc-s} gives an error that grows with
  the converted power. It is a common failure mode.
\item \textbf{Signs.} The theorem fixes the sign of a scattered
  amplitude but not its magnitude. For example it fixes the sign of the
  toroidal coefficient in \texttt{FirstOrderContrastOperator.tex}.
  A family that conserves energy on its own passes the test whether it
  is present or not.
\item \textbf{Approximate T-matrices.} A Born or first-order T-matrix
  is real to leading order and has no radiation damping. Its forward
  amplitude has an imaginary part only at second order in contrast, so
  it cannot satisfy the theorem exactly
  \citep{GubernatisDomanyKrumhansl1977}. The residual measures the
  missing self-consistency and is not a coding error.
\end{itemize}

% =============================================================================
\section{Conventions and checks in this repository}
\label{sec:sign}
% =============================================================================

Equations \eqref{eq:ot-p} and \eqref{eq:ot-s} assume $e^{-i\omega t}$
and far-field amplitudes defined by \eqref{eq:farfield}. Under
$e^{+i\omega t}$ the imaginary part changes sign, and so does any
convention that puts an extra minus sign on the far-field amplitude. Every
checker in \texttt{scattered\_field.py} now uses the standard sign,
$\sext = +(4\pi/k)\Imag f(\bd)$, as does
\texttt{docs/cube\_tmatrix\_closedform.tex}. An earlier minus sign there
compensated a force-term sign error in the cube far field, and the two
were removed together.

For P incidence on the lossless $T_{27}$ cube, $\sext/\ssc$ is
$0.99997$ at $ka = 0.05$ and $0.9986$ at $ka = 0.3$. For a pure-modulus
contrast it is $1.00002$. Getting there required the modulus radiation
reaction to be treated like the density one:
\begin{itemize}
\item \textbf{Kept in the far field.} At forward incidence the incident
  strain $i\kp\bd\bd$ is imaginary, so $\Imag\Delta\mathbf{c}^*$ is
  exactly what enters $\Imag f_P(\bd)$. A far field built from
  $\Real\Delta\mathbf{c}^*$ alone gave a pure-modulus cube $\sext = 0$
  and the full cube a ratio of $0.61$.
\item \textbf{Built from the effective contrast.} The reaction is the
  radiated power of the dipole the far field actually emits,
  $\Imag\Delta\mathbf{c}^* = -(2V/\omega)\,
  \Real\Delta\mathbf{c}^*:\mathbf{K}:\Real\Delta\mathbf{c}^*$, with
  $\mathbf{K}$ the isotropic radiated-power kernel. Built from the bare
  contrast it overshoots, giving a ratio near $1.02$. The exact Mie
  sphere obeys this identity with its own effective contrasts to
  $10^{-4}$.
\end{itemize}
The closed form and its checks are in
\texttt{docs/cube\_tmatrix\_closedform.tex} (eq.~modulus-reaction).
This case belongs with the \emph{signs} and \emph{mode-conversion}
entries of \S\ref{sec:check}: a missing radiation reaction is a
bookkeeping defect, and the theorem detects it exactly.

\bibliography{psvsh_optical_theorem}

\end{document}
